\clearpage
\section*{Literatur- und Quellenverzeichnis}
\addcontentsline{toc}{section}{Literatur- und Quellenverzeichnis}

\begin{description}
  \item[LaTeX-Vorlage] Stefan Macke: LaTeX-Vorlage zur IHK-Projektdokumentation
  für Fachinformatiker Anwendungsentwicklung, lizenziert unter Creative Commons
  Namensnennung - Weitergabe unter gleichen Bedingungen 4.0 International
  (CC BY-SA 4.0).\\
  \url{http://fiae.link/LaTeXVorlageFIAE}\\
  \url{http://creativecommons.org/licenses/by-sa/4.0/}

  \item[Python] Python Software Foundation: Python~3.12 Documentation.\\
  \url{https://docs.python.org/3.12/}

  \item[FastAPI] FastAPI: Offizielle Dokumentation zum Webframework, Routing,
  Dependency Injection und API-Aufbau.\\
  \url{https://fastapi.tiangolo.com/}

  \item[SQLModel] SQLModel: Offizielle Dokumentation zur Modellierung und
  Nutzung relationaler Datenbanken in Python.\\
  \url{https://sqlmodel.tiangolo.com/}

  \item[PlantUML] PlantUML: Dokumentation zur Erstellung von UML-Diagrammen.\\
  \url{https://plantuml.com/de/}

  \item[Argon2] argon2-cffi: Dokumentation zu Passwort-Hashing mit
  \Code{PasswordHasher}.\\
  \url{https://argon2-cffi.readthedocs.io/en/stable/howto.html}

  \item[W3Schools] W3Schools: HTML Forms, Grundlagen zu Formularen und
  Eingabefeldern.\\
  \url{https://www.w3schools.com/html/html_forms.asp}

  \item[Stack Overflow] Stack Overflow: Diskussion zu FastAPI OAuth2 Scope und
  rollenbasierter Zugriffskontrolle.\\
  \url{https://stackoverflow.com/questions/64812279/fastapi-oauth2-scope}

  \item[Reddit] Reddit, r/FastAPI: Diskussion zu FastAPI, SQLModel und
  API-Aufbau im praktischen Einsatz.\\
  \url{https://www.reddit.com/r/FastAPI/comments/pgdd6p}

  \item[KI-Unterstützung] OpenAI: ChatGPT, genutzt zur Informationsbeschaffung, sprachlichen
  Überarbeitung, Formulierungshilfe und CSS-Angaben.\\
  \url{https://openai.com/chatgpt/overview/}
\end{description}
